\colhead{对应知识清单}{第1—10题}
\field{第1题}{元素与集合用 $\in$；集合与集合用 $\subseteq$。交集取公共元素，并集取所有元素，补集必须相对于全集。}
\field{第2题}{区间的开闭端点表示能否取到。并集恰好等于给定区间，要同时做到无遗漏、无超出；连接点只需由一个集合覆盖。}
\field{第3题}{$p\Rightarrow q$ 时p是q的充分条件，q是p的必要条件；两方向均能推出才是充要。否定某个推出关系，只需一个反例。}
\field{第4题}{全称命题的否定改用存在量词，存在命题的否定改用全称量词；量词范围保持不变。$>$ 的否定是 $\le$，$<$ 的否定是 $\ge$。}
\field{第5题}{同加减保序；乘除正数保序，乘除负数反向。非负不等于正，$d^2$ 在 $d=0$ 时为0。}
\field{第6题}{先统一为正的二次项系数，再找根、画数轴、判端点。含参根须比较顺序；重根时回到原不等式判断。}
\field{第7题}{对非负数u、v，有 $u+v\ge2\sqrt{uv}$，等号当且仅当 $u=v$。求最值必须检查正负条件、定值结构和等号可取性。}
\field{第8题}{对所有x都非负，等价于可取得的最小值非负；对所有x都为正，须结合最小值是否取得判断。二次式配方后可直接找最值。}
\field{第9题}{函数要求定义域中每个数都有唯一像；允许多个自变量有相同函数值。相同函数须定义域、对应关系都相同，值域是实际取到的像的集合。}
\field{第10题}{偶次根式的被开方数须非负，分母须非零；多个限制同时满足，取交集。根号内等于0时，还要检查分母。}
\newpage
\colhead{对应知识清单}{第11—20题}
\field{第11题}{求 $f(f(a))$ 先求内层，再把结果作为新输入；分段函数每次都要重新选段。已知 $h(x+c)$ 时，先看括号内整体。}
\field{第12题}{定义域看图象覆盖的横坐标，值域看纵坐标；读函数值先找横坐标。图象解不等式时，要保留水平线交点并结合原定义域。}
\field{第13题}{二次函数先配方找对称轴，再与定义域比较。闭区间最值须检查顶点是否在内以及两端点；求出所有取等位置，值域写成集合。}
\field{第14题}{定义证明须任取区间内两个有大小顺序的数，作差、变形、定号并下结论。二次函数区间单调问题取决于对称轴的位置。}
\field{第15题}{奇函数定义域关于原点对称，满足 $f(-x)=-f(x)$；若0在定义域中，则 $f(0)=0$。已知半轴解析式时先把输入变到已知半轴。}
\field{第16题}{$g(x)+g(-x)$ 的定义域需使两个输入均合法；在对称域内它是偶函数。利用轴对称补出另一半图象及全部最值点。}
\field{第17题}{正底数有 $a^{m/n}=(\sqrt[n]a)^m$，$a^{-r}=1/a^r$。$\sqrt{u^2}=|u|$；奇次实根保留被开方数的符号。}
\field{第18题}{幂函数由给定点确定指数，再按原式求定义域。奇偶性先查对称域；单调区间不能跨越无定义点。负分数指数还须排除根式为0。}
\field{第19题}{指数函数底数须大于0且不为1；底数大于1时保序，介于0与1时反序。带绝对值的指数先找分界点，再看内外两层的增减。}
\field{第20题}{指数模型必须保留实际定义域。衰减模型底数介于0与1之间；同底指数比较反向；趋近0与实际等于0不同。}
